\FloatBarrier

\ChapterImageStar[cap:pmv]{Evaluación de la notación NDTI}{./images/fondo.png}\label{cap:pmv}
\mbox{}\\

\subsection{Enfoque metodológico}
\noindent
La evaluación de la notación NDTI se desarrolló mediante un enfoque mixto con un diseño comparativo entre grupos, centrado principalmente en la aplicación práctica de la notación a través de actividades de diagramación.

Para ello, se conformaron dos grupos de participantes, cada uno compuesto por cinco estudiantes. Al primer grupo se le presentó previamente la notación NDTI, incluyendo sus elementos, capas y conectores, mientras que el segundo grupo desarrolló la actividad de diagramación de manera libre, sin el uso de una notación específica. Esta configuración permitió establecer una comparación directa entre los modelos generados, con el fin de analizar el impacto del uso de la notación en la calidad, claridad y coherencia de los diagramas resultantes.

El componente principal de la evaluación se fundamenta en el análisis de los diagramas elaborados por ambos grupos, permitiendo identificar diferencias en la selección de elementos, el uso de relaciones y la estructuración de los modelos. De manera complementaria, se aplicó una encuesta estructurada tipo Likert únicamente al grupo que utilizó la notación NDTI, con el propósito de recoger su percepción respecto a la claridad, facilidad de uso y utilidad de la misma.

Asimismo, se incluyeron preguntas abiertas que permitieron identificar dificultades específicas, patrones de uso y oportunidades de mejora en la notación. De esta manera, la evaluación integra tanto la comparación entre resultados de modelado como la percepción de los usuarios, proporcionando una base más robusta para el análisis de la efectividad de la notación propuesta.

\subsection{Diseño del ejercicio de evaluación}
\noindent
La evaluación se estructuró en torno a una única actividad de diagramación, en la cual los participantes debían representar un escenario mediante el uso de la notación NDTI.

El ejercicio propuesto consistió en un enunciado de carácter transversal que describe un sistema de información involucrando múltiples niveles de la arquitectura de TI. En este sentido, el problema planteado requiere la integración de elementos correspondientes a distintas capas de la notación, tales como la capa física, redes y comunicaciones, tecnológica, de aplicación y de negocio.

El diseño de este ejercicio tuvo como objetivo evaluar la capacidad de los participantes para construir un modelo coherente que articule correctamente los distintos componentes del sistema, así como el uso adecuado de los conectores transversales definidos en la notación.

De esta manera, al emplear un único ejercicio multicapa, se garantiza la comparabilidad entre los resultados obtenidos por los distintos grupos, permitiendo analizar de forma consistente el impacto del uso de la notación en la calidad y estructura de los diagramas generados.

\subsection{Instrumentos de recolección de información}
\noindent
La recolección de información se realizó mediante dos instrumentos complementarios, los cuales aportan perspectivas diferentes pero igualmente relevantes para la evaluación de la notación NDTI:

\begin{itemize}
\item \textbf{Diagramas elaborados por los participantes:} Permiten analizar la aplicación práctica de la notación en la resolución del ejercicio propuesto. A través de estos se evalúan aspectos como la selección de elementos, el uso de relaciones, la coherencia estructural del modelo y la correspondencia con el enunciado.

\item \textbf{Encuesta de percepción:} Instrumento aplicado al grupo que utilizó la notación NDTI, diseñado para recopilar la experiencia de los participantes durante el ejercicio. La encuesta incluyó preguntas cerradas en escala Likert y preguntas abiertas, orientadas a identificar el nivel de comprensión, facilidad de uso, utilidad percibida y dificultades encontradas.
\end{itemize}

\noindent
Ambos instrumentos se consideran fundamentales dentro del proceso de evaluación, ya que permiten contrastar la evidencia empírica derivada del desempeño en la actividad de diagramación con la percepción de los participantes sobre el uso de la notación.

\subsection{Procedimiento}
\noindent
El proceso de evaluación se desarrolló en las siguientes etapas:

\begin{enumerate}
\item \textbf{Conformación de los grupos:} Se organizaron dos grupos de participantes, cada uno compuesto por cinco estudiantes. Uno de los grupos fue asignado como grupo experimental (uso de la notación NDTI) y el otro como grupo de referencia (diagramación libre).

\item \textbf{Introducción a la notación (grupo experimental):} Al grupo experimental se le presentó la notación NDTI, incluyendo sus elementos, capas y conectores, junto con ejemplos de referencia que facilitaran su comprensión. El grupo de referencia no recibió esta instrucción.

\item \textbf{Desarrollo de la actividad de diagramación:} Ambos grupos resolvieron el mismo ejercicio multicapa, construyendo un diagrama que representara el escenario propuesto. El grupo experimental debía utilizar la notación NDTI, mientras que el grupo de referencia podía representar el problema de manera libre.

\item \textbf{Aplicación de la encuesta:} Una vez finalizada la actividad, se aplicó la encuesta de evaluación únicamente al grupo experimental, con el fin de recoger su percepción sobre el uso de la notación.

\item \textbf{Recolección de resultados:} Se recopilaron los diagramas elaborados por ambos grupos y las respuestas de la encuesta, constituyendo la base para el análisis comparativo de los resultados.
\end{enumerate}

\subsection{Criterios de evaluación}
\noindent
El análisis de los resultados se estructuró a partir de dos dimensiones complementarias: (i) el desempeño en la actividad de diagramación y (ii) la percepción de los participantes sobre la notación NDTI.

\subsubsection{Evaluación de los diagramas}
\noindent
La evaluación de los diagramas se realizó de forma comparativa entre ambos grupos, considerando criterios aplicables independientemente del uso de una notación específica:

\begin{itemize}
\item \textbf{Correspondencia con el enunciado:} Grado en que el modelo representa correctamente los elementos, relaciones y comportamientos descritos en el problema planteado.

\item \textbf{Estructuración del modelo:} Nivel de organización, claridad visual y disposición lógica de los elementos dentro del diagrama.

\item \textbf{Representación de relaciones:} Coherencia en la forma en que se modelan interacciones, dependencias y flujos entre los distintos componentes del sistema.

\item \textbf{Completitud del modelo:} Nivel de cobertura del escenario descrito, considerando la inclusión de los componentes relevantes y sus interacciones.
\end{itemize}

\noindent
Adicionalmente, para el grupo experimental que utilizó la notación NDTI, se incorporaron los siguientes criterios específicos:

\begin{itemize}
\item \textbf{Uso adecuado de elementos NDTI:} Correcta selección y aplicación de los elementos de la notación según su definición semántica.

\item \textbf{Uso adecuado de conectores:} Aplicación coherente de los conectores transversales (asociación, flujo, dependencia, realización, trigger, especialización y comunicación) de acuerdo con su propósito dentro del modelo.
\end{itemize}

\subsubsection{Evaluación de la percepción (encuesta)}
\noindent
La información obtenida mediante la encuesta fue analizada a partir de cuatro ejes principales, los cuales permiten interpretar la experiencia de uso de la notación:

\begin{itemize}
\item \textbf{Comprensión de la notación:} Nivel en que los participantes lograron entender los elementos, capas y relaciones propuestas.

\item \textbf{Uso de los elementos:} Facilidad percibida para seleccionar y aplicar correctamente los elementos dentro del diagrama.

\item \textbf{Expresividad y utilidad:} Capacidad de la notación para representar el problema de manera clara, completa y adecuada al contexto planteado.

\item \textbf{Experiencia general:} Valoración global del uso de la notación, incluyendo percepción de dificultad, claridad y aplicabilidad.
\end{itemize}

\noindent
La integración de ambas dimensiones permitió contrastar el desempeño observado en los diagramas con la percepción de los participantes, facilitando una evaluación más completa de la efectividad de la notación NDTI.

\subsection{Método de análisis}
\noindent
El análisis de los datos se realizó mediante un enfoque mixto, integrando técnicas cualitativas y cuantitativas con el fin de evaluar tanto el desempeño en la actividad de diagramación como la percepción de los participantes sobre la notación NDTI.

\begin{itemize}
\item \textbf{Análisis de diagramas:} Los diagramas elaborados por ambos grupos fueron evaluados de manera comparativa a partir de los criterios definidos previamente. Este análisis combinó una valoración estructurada de los criterios (correspondencia, estructuración, relaciones y completitud) con un análisis cualitativo orientado a identificar patrones de representación, diferencias entre grupos, errores recurrentes y dificultades en la modelación del problema.

\item \textbf{Análisis de la encuesta:} Las respuestas de la encuesta aplicada al grupo experimental fueron analizadas mediante estadística descriptiva, calculando promedios por pregunta y por cada uno de los ejes de evaluación (comprensión de la notación, uso de los elementos, expresividad y utilidad, y experiencia general). Esto permitió identificar tendencias en la percepción de los participantes.

\item \textbf{Análisis de respuestas abiertas:} Las respuestas abiertas fueron clasificadas y categorizadas, permitiendo complementar la interpretación de los resultados cuantitativos. Este análisis facilitó la identificación de dificultades específicas, percepciones recurrentes y posibles oportunidades de mejora en la notación.

\end{itemize}

\noindent
Finalmente, los resultados obtenidos en ambas dimensiones fueron integrados con el fin de contrastar el desempeño observado en los diagramas con la percepción de los participantes, permitiendo una evaluación más completa de la efectividad de la notación NDTI.

